\subsection{Componentwise control up the cyclotomic tower}
\label{sec:componentwise-iwasawa-control}

The finite class-component vanishings also remove specified odd torsion
companions in the cubic criterion. This conclusion concerns entire
Iwasawa components, not just one specialization of their characteristic
power series.

\begin{proposition}[K17-COMPONENT-CONTROL]
\label{prop:componentwise-iwasawa-control}
Let $p$ be odd, put
\[
 F_n=\mathbf Q(\mu_{p^{n+1}}),\qquad
 A_n=\operatorname{Cl}(F_n)\{p\},\qquad
 X=\varprojlim_n A_n,
\]
where the transition maps are norms. Let
$\Gamma=\operatorname{Gal}(F_\infty/F_0)$,
$\Delta=\operatorname{Gal}(F_0/\mathbf Q)$, and
$\Lambda=\mathbf Z_p[[\Gamma]]$. Use the Teichm\"uller splitting
to let $\Delta$ act on every layer. For every character $\theta$ of
$\Delta$ one has
\[
 A_0^{(\theta)}=0
 \quad\Longrightarrow\quad
 A_n^{(\theta)}=0\ \hbox{for every }n\geq0
 \quad\Longrightarrow\quad X^{(\theta)}=0.
 \tag{CC.1}
\]
For $p\geq5$, let $M$ be the Galois group of the maximal abelian
pro-$p$ extension of $F_\infty$ unramified outside $p$.
If $\theta$ is even and $\theta\ne1,\omega^2$, then
\[
 A_0^{(\theta)}=0
 \quad\Longrightarrow\quad
 M^{(\omega\theta^{-1})}\text{ is free of rank one over }\Lambda.
 \tag{CC.2}
\]
Here $\omega$ is the Teichm\"uller character. The freeness in
(CC.2) is integral.
\end{proposition}

\begin{proof}
For each $n$, the extension $F_n/F_0$ is cyclic of degree $p^n$.
Its only ramified prime is the unique prime above $p$, and that
prime is totally ramified. Write $G_n=\operatorname{Gal}(F_n/F_0)$.
The norm map gives an isomorphism
\[
                 (A_n)_{G_n}\ \xrightarrow{\ \sim\ }\ A_0.
 \tag{CC.3}
\]
This is precisely the combination of the norm surjectivity in
\cite[Corollary~1.3.9, p.~17]{SharifiIwasawaFull2026} and the
one-prime coinvariant injection in
\cite[Corollary~1.3.12, p.~18]{SharifiIwasawaFull2026}.
The infinite-level formulation is also stated explicitly in
\cite[Proposition~3.2.11, equation~(3.2.1), pp.~80--81;
Remark~3.2.12, p.~81]{SharifiIwasawaFull2026}.

Here is the class-field argument for (CC.3). Let $L$ be the largest
subfield of the $p$-Hilbert class field of $F_n$ which is abelian
over $F_0$. Since $G_n$ is cyclic, the commutator subgroup of the
Galois group over $F_0$ is generated by the commutators with its
abelian normal subgroup over $F_n$. Consequently
$\operatorname{Gal}(L/F_n)=(A_n)_{G_n}$.
In the abelian group $\operatorname{Gal}(L/F_0)$, the inertia
group $J$ at the unique prime above $p$ maps isomorphically onto
$G_n$: its kernel is trivial because $L/F_n$ is unramified, and
its image is all of $G_n$ by total ramification. Thus
$\operatorname{Gal}(L/F_0)$ is the direct product of $J$ and
$\operatorname{Gal}(L/F_n)$. The field $L^J$ is unramified over
$F_0$ at every place and is abelian of $p$-power degree. Conversely,
the $p$-Hilbert class field of $F_0$, after composition with $F_n$,
lies in $L$. Its intersection with $F_n$ is $F_0$, again by total
ramification. Hence $L^J$ is that Hilbert class field, proving
(CC.3) under the Artin identification. The map is the norm map,
so it is canonical.

Norms commute with every automorphism in $\Delta$, and the
Teichm\"uller subgroup commutes with $G_n$. The integral
idempotents of $\mathbf Z_p[\Delta]$ therefore give
\[
 (A_n^{(\theta)})_{G_n}\simeq A_0^{(\theta)}.
\]
Suppose the right-hand side is zero and put
$V=A_n^{(\theta)}/pA_n^{(\theta)}$. If $\sigma$ generates $G_n$,
then $V=(\sigma-1)V$, while
$(\sigma-1)^{p^n}=\sigma^{p^n}-1=0$ on $V$.
Iteration gives $V=0$. Since $A_n^{(\theta)}$ is a finite
$p$-group, $A_n^{(\theta)}=pA_n^{(\theta)}$ implies
$A_n^{(\theta)}=0$. This proves the first implication of (CC.1)
at every finite layer. Taking the norm inverse limit proves the
second implication; no exactness of an inverse limit is needed.
In particular there is no mixing of distinct $\Delta$-characters.

For (CC.2), put $\chi=\omega^{-1}\theta$, so
$\chi^{-1}=\omega\theta^{-1}$. At tame level one, the assumptions
of Wake--Wang-Erickson's setup hold: $\theta$ is even and primitive
because it is nontrivial modulo $p$; the exceptional even character
$\omega^2$ is excluded. Their full duality calculation gives
\[
 \operatorname{Ext}^1_\Lambda(M_{\chi^{-1}}^{\#}(1),\Lambda)
       =X_\theta,
 \qquad
 \operatorname{Ext}^2_\Lambda(M_{\chi^{-1}}^{\#}(1),\Lambda)=0,
\]
and proves that $X_\theta=0$ is equivalent to
$M_{\chi^{-1}}^{\#}(1)$ being free of rank one
\cite[\S2.1, pp.~2--3; Proposition~2.2.6 and its full proof,
p.~4]{WakeWangErickson2017}. The module on the left is their
$p$-ramified, fraktur module; the module $X$ on the right is
unramified. By (CC.1), the right-hand side is zero.
Inversion and the Tate twist transport the Iwasawa algebra by
automorphisms, so undoing them preserves integral rank-one
freeness. This proves (CC.2).
\end{proof}

\begin{corollary}[K17-ODD-COMPANION-EXCLUSION]
\label{cor:odd-companion-exclusion}
Under the stipulated arithmetic theorem and the conclusions proved
above, for every odd prime $p$ one has
\[
 A_n^{[-2j]}=0\quad(n\geq0),\qquad X^{[-2j]}=0,
                  \qquad 0\leq j\leq9.
 \tag{CC.4}
\]
For $p\geq5$, the actual $p$-ramified component $M^{[2j+1]}$ is
integrally free of rank one for each $0\leq j\leq9$.
Consequently the odd torsion submodule in
Theorem~\ref{thm:odd-torsion-cubic} has zero components in all
these residue characters. In particular no Triple-NV condition
is required for an odd companion in one of them.
\end{corollary}

\begin{proof}
For $1\leq j\leq9$, Corollary~\ref{cor:nine-degrees-period}
and the exact vanishing comparison (K2) give $A_0^{[-2j]}=0$.
For $j=0$, the trivial class component is zero because norm to
$\mathbf Q$ followed by extension is multiplication by $p-1$
on that component, and the class group of $\mathbf Q$ is zero.
Apply (CC.1) to obtain (CC.4). For a nonexceptional character,
(CC.2) applies with
$\theta=\omega^{-2j}$ and
$\omega\theta^{-1}=\omega^{2j+1}$.

For completeness, the excluded odd characters $\omega$ and
$\omega^{-1}$ already have integral rank-one freeness without
Greenberg. The previously established Coleman calculation uses
\[
0\longrightarrow\mathbf Z_p(1)\longrightarrow U_{\rm loc}
 \longrightarrow\mathbf Z_p[[\Delta\times\Gamma]]
 \longrightarrow\mathbf Z_p(1)\longrightarrow0
\]
and $X_\omega=X_{\omega^{-1}}=0$. Global minus units contribute
only the root-of-unity term. Thus $M_{\omega^{-1}}\simeq\Lambda$,
whereas
$M_\omega\simeq\ker(\Lambda\to\mathbf Z_p(1))
=(\gamma-u)\Lambda$, with $u=\chi_p(\gamma)$.
Both are free of rank one. The exact source locators are
\cite[Theorem~5.5.31, p.~148; Proposition~6.1.1,
pp.~151--152]{SharifiIwasawaNotes} and
\cite[proof of Lemma~5.2.7, p.~13]{WakeWangErickson2017}.
This also covers repetitions or exceptional residue values in
the finite list. A free $\Lambda$-module has no $\Lambda$-torsion,
which proves the last assertion.
\end{proof}

The same argument works for any further set of proved vanishings
$K_{4j}(\mathbf Z)\{p\}=0$: each such $j$ removes the entire
even component $X^{[-2j]}$ and the matching odd torsion companion
in character $[2j+1]$. For a fixed large prime, the finite list
above does not exhaust all even characters. It therefore neither
proves Greenberg nor supplies the missing cubic nonvanishing in
the remaining characters. It gives an explicit restriction on
where that obstruction can still occur.
